% Generated draft: requires mathematical and source review.
\documentclass[11pt]{article}
\usepackage[margin=25mm]{geometry}
\usepackage{fontspec}
\setmainfont{Latin Modern Roman}
\usepackage{amsmath,amssymb,mathtools,unicode-math}
\setmathfont{Latin Modern Math}
\usepackage{enumitem,needspace,xurl,hyperref}
\setlength{\parindent}{0pt}
\setlength{\parskip}{5pt}
\setlength{\emergencystretch}{3em}
\newcommand{\N}{\mathbb{N}}
\newcommand{\Z}{\mathbb{Z}}
\newcommand{\Q}{\mathbb{Q}}
\newcommand{\R}{\mathbb{R}}
\newcommand{\C}{\mathbb{C}}
\newcommand{\F}{\mathbb{F}}
\newcommand{\A}{\mathbb{A}}
\newcommand{\PP}{\mathbb{P}}
\newcommand{\E}{\mathbb{E}}
\newcommand{\eps}{\varepsilon}
\newcommand{\dd}{\,\mathrm{d}}
\begin{document}
\section*{UnsolvedMath ID: 30000420}
\textbf{Source problem code:} catalogue entry
\subsection{Definitions and mathematical statement}
An orientation of a graph assigns one of the two directions to each edge. Globally consistent local orientation data mean, at minimum, that the endpoints of every edge agree on its direction; this does not by itself imply acyclicity, a global vertex order, or bounded outdegree. The source's more specific orientation convention, if any, is not reproduced.

A proper vertex coloring is a map \(c:V(G)\to\{1,\ldots,q\}\) satisfying \(c(u)\ne c(v)\) for every edge \(uv\). The source asks whether providing the stated orientations in a distributed graph-coloring model permits faster coloring algorithms. To make this a precise comparison one must fix a graph class, allowed palette \(q\), communication and computational model, available initial information, and a resource such as the number of communication rounds. One could then compare the optimal resource bound with the orientation supplied against the bound in the same model without it. No particular asymptotic speedup is asserted in the supplied record.

The extract does not specify deterministic versus randomized algorithms, vertex identifiers or anonymity, message-size limits, synchrony, degree bounds, initial coloring, failure probability, or whether the orientation is free input or must be computed at a cost. These choices can change both the baseline and the benefit. Thus the question is an algorithmic research direction, not a theorem about an invented LOCAL or CONGEST model or a specified color bound.

The supplied assessment checked 2026-08-21 expressly says that the formulation is too incomplete to identify a solved theorem and assigns open status under its curation policy. That status is not evidence of a particular proved lower bound or a fully identified unresolved complexity comparison.
\subsection{Short English statement}
Can consistent edge directions speed up distributed proper coloring? The coloring budget, communication model, orientation guarantees, and baseline running time must be specified before an asymptotic algorithmic claim can be stated.
\subsection{Sources}
\begin{itemize}
\item Local UnsolvedMath registry, permanent ID 30000420, catalogue entry, \emph{Faster Distributed Coloring from Consistent Edge Orientations}.
\item \emph{Algorithmic Graph Theory} (2006), Oberwolfach Reports, pp. 379--460, DOI \url{https://doi.org/10.4171/owr/2006/07}.
\item Supplied literature triage checked 2026-08-21 cites only that report; no complete model definition, algorithm, or lower-bound theorem is supplied.
\end{itemize}
\end{document}
